% Generated draft: requires mathematical and source review.
\documentclass[11pt]{article}
\usepackage[margin=25mm]{geometry}
\usepackage{fontspec}
\setmainfont{Latin Modern Roman}
\usepackage{amsmath,amssymb,mathtools,unicode-math}
\setmathfont{Latin Modern Math}
\usepackage{enumitem,needspace,xurl,hyperref}
\setlength{\parindent}{0pt}
\setlength{\parskip}{5pt}
\setlength{\emergencystretch}{3em}
\newcommand{\N}{\mathbb{N}}
\newcommand{\Z}{\mathbb{Z}}
\newcommand{\Q}{\mathbb{Q}}
\newcommand{\R}{\mathbb{R}}
\newcommand{\C}{\mathbb{C}}
\newcommand{\F}{\mathbb{F}}
\newcommand{\A}{\mathbb{A}}
\newcommand{\PP}{\mathbb{P}}
\newcommand{\E}{\mathbb{E}}
\newcommand{\eps}{\varepsilon}
\newcommand{\dd}{\,\mathrm{d}}
\begin{document}
\section*{UnsolvedMath ID: 30000270}
\textbf{Source problem code:} catalogue entry
\subsection{Definitions and mathematical statement}
For a contact manifold \((V,\alpha)\), its standard symplectization is \((\R\times V,\dd(e^a\alpha))\), with \(a\) the real coordinate and an adapted almost-complex structure. The source asks why a harmonic one-form is needed in an intended perturbation of the pseudoholomorphic-curve equation for a Riemann surface \((\Sigma,j)\) other than \(S^2\). It does not reproduce that equation, the contact data, or the role assigned to the perturbing form.

A real one-form \(\gamma\) on a Riemann surface is harmonic, in the ordinary smooth compact setting, when
\[
\dd\gamma=0,\qquad\dd^*\gamma=0,
\]
for a compatible metric; equivalently one can express the second condition using the complex structure and the Hodge star. On a closed surface, harmonic forms represent real first cohomology. In particular \(H^1(S^2;\R)=0\), while a closed genus-\(g\) surface has a \(2g\)-dimensional space of harmonic real one-forms. These facts explain a possible distinction between the sphere and positive genus, but they do not by themselves prove necessity in the missing nonlinear equation.

For example, if a source equation requires a closed one-form \(\beta\) to be represented by \(\dd a\), global solvability requires its periods to vanish; allowing \(\beta=\dd a+\gamma\) can retain the harmonic cohomology component. This is a conditional description of a possible period obstruction, not a reconstruction of the absent perturbed equation or a claim that its particular form is always closed.

The source question is to identify precisely what obstruction, existence statement, or geometric requirement makes the harmonic term necessary, and under what hypotheses it could be omitted. It does not assert that every solution with nonspherical domain has a nonzero harmonic term. Punctures, ends, energy, and asymptotic conditions would change the relevant harmonic-form space and must be specified; the closed-surface statements above are not applied to punctured domains without them.

The supplied assessment checked 2026-08-21 reports that no general theorem eliminating or explaining the necessity for the stated equation was located. Because the equation and the intended sense of necessity are missing, this remains a conceptual source-dependent question rather than one fully quantified theorem.
\subsection{Short English statement}
Explain the role of the harmonic one-form in the intended perturbed pseudoholomorphic equation on nonspherical domains. Nontrivial periods may be relevant, but the equation and the precise meaning of necessity must be recovered before that explanation becomes a theorem.
\subsection{Sources}
\begin{itemize}
\item Local UnsolvedMath registry, permanent ID 30000270, catalogue entry, \emph{Harmonic One-Forms in Perturbed Pseudoholomorphic-Curve Equations}.
\item \emph{Dynamical Systems} (2006), Oberwolfach Reports, pp. 1743--1798, DOI \url{https://doi.org/10.4171/owr/2005/31}.
\item Supplied literature triage checked 2026-08-21 cites only the original report; no explicit perturbed equation or separate theorem citation is supplied.
\end{itemize}
\end{document}
